\begin{lemma}
\label{lemma-perfect-modulo-two-ideals}
Let $R$ be a ring. Let $I, J \subset R$ be ideals.
Let $K$ be an object of $D(R)$. Assume that
\begin{enumerate}
\item $K \otimes_R^\mathbf{L} R/I$ is perfect in $D(R/I)$, and
\item $K \otimes_R^\mathbf{L} R/J$ is perfect in $D(R/J)$.
\end{enumerate}
Then $K \otimes_R^\mathbf{L} R/IJ$ is perfect in $D(R/IJ)$.
\end{lemma}

\begin{proof}
It is clear that we may assume replace $R$ by $R/IJ$ and $K$ by
$K \otimes_R^\mathbf{L} R/IJ$. Then $R \to R/(I \cap J)$ is
a surjection whose kernel has square zero. Hence by
Lemma \ref{lemma-perfect-modulo-nilpotent-ideal}
it suffices to prove that $K \otimes_R^\mathbf{L} R/(I \cap J)$ is
perfect. Thus we may assume that $I \cap J = 0$.

\medskip\noindent
We prove the lemma in case $I \cap J = 0$. First, we may represent $K$
by a K-flat complex $K^\bullet$ with all $K^n$ flat, see
Lemma \ref{lemma-K-flat-resolution}. Then we see that we have a short
exact sequence of complexes
$$
0 \to
K^\bullet \to K^\bullet/IK^\bullet \oplus K^\bullet/JK^\bullet \to
K^\bullet/(I + J)K^\bullet \to 0
$$
Note that $K^\bullet/IK^\bullet$ represents $K \otimes^\mathbf{L}_R R/I$
by construction of the derived tensor product. Similarly for
$K^\bullet/JK^\bullet$ and $K^\bullet/(I + J)K^\bullet$.
Note that $K^\bullet/(I + J)K^\bullet$ is a perfect complex
of $R/(I + J)$-modules, see
Lemma \ref{lemma-pull-perfect}.
Hence the complexes $K^\bullet/IK^\bullet$, and
$K^\bullet/JK^\bullet$ and $K^\bullet/(I + J)K^\bullet$
have finitely many nonzero cohomology groups
(since a perfect complex has finite Tor-amplitude, see
Lemma \ref{lemma-perfect}). We conclude that $K \in D^b(R)$ by the
long exact cohomology sequence associated to short exact sequence
of complexes displayed above. In particular we assume $K^\bullet$
is a bounded above complex of free $R$-modules (see
Derived Categories, Lemma \ref{derived-lemma-subcategory-left-resolution}).

\medskip\noindent
We will now show that $K$ is perfect using the criterion of
Proposition \ref{proposition-perfect-is-compact}. Thus we let
$E_j \in D(R)$ be a family of objects parametrized by a set $J$.
We choose complexes $E_j^\bullet$ with flat terms
representing $E_j$, see for example Lemma \ref{lemma-K-flat-resolution}.
It is clear that
$$
0 \to
E_j^\bullet \to
E_j^\bullet/IE_j^\bullet \oplus E_j^\bullet/JE_j^\bullet \to
E_j^\bullet/(I + J)E_j^\bullet \to 0
$$
is a short exact sequence of complexes. Taking direct sums we obtain a
similar short exact sequence
$$
0 \to
\bigoplus E_j^\bullet \to
\bigoplus E_j^\bullet/IE_j^\bullet \oplus E_j^\bullet/JE_j^\bullet \to
\bigoplus E_j^\bullet/(I + J)E_j^\bullet \to 0
$$
(Note that $- \otimes_R R/I$ commutes with direct sums.)
This short exact sequence determines a distinguished triangle in $D(R)$, see
Derived Categories, Lemma \ref{derived-lemma-derived-canonical-delta-functor}.
Apply the homological functor $\Hom_{D(R)}(K, -)$ (see
Derived Categories, Lemma \ref{derived-lemma-representable-homological})
to get a commutative diagram
$$
\xymatrix{
\bigoplus \Hom_{D(R)}(K^\bullet, E_j^\bullet/(I + J))[-1] \ar[r] \ar[d] &
\Hom_{D(R)}(K^\bullet, \bigoplus E_j^\bullet/(I + J))[-1] \ar[d] \\
\bigoplus \Hom_{D(R)}(K^\bullet, E_j^\bullet/I \oplus E_j^\bullet/J)[-1]
\ar[r] \ar[d] &
\Hom_{D(R)}(K^\bullet, \bigoplus E_j^\bullet/I \oplus E_j^\bullet/J)[-1]
\ar[d] \\
\bigoplus \Hom_{D(R)}(K^\bullet, E_j^\bullet) \ar[r] \ar[d] &
\Hom_{D(R)}(K^\bullet, \bigoplus E_j^\bullet) \ar[d] \\
\bigoplus \Hom_{D(R)}(K^\bullet, E_j^\bullet/I \oplus E_j^\bullet/J)
\ar[r] \ar[d] &
\Hom_{D(R)}(K^\bullet, \bigoplus E_j^\bullet/I \oplus E_j^\bullet/J)
\ar[d] \\
\bigoplus \Hom_{D(R)}(K^\bullet, E_j^\bullet/(I + J)) \ar[r] &
\Hom_{D(R)}(K^\bullet, \bigoplus E_j^\bullet/(I + J))
}
$$
with exact columns. It is clear that, for any complex $E^\bullet$
of $R$-modules we have
\begin{align*}
\Hom_{D(R)}(K^\bullet, E^\bullet/I) & =
\Hom_{K(R)}(K^\bullet, E^\bullet/I) \\
& =
\Hom_{K(R/I)}(K^\bullet/IK^\bullet, E^\bullet/I) \\
& =
\Hom_{D(R/I)}(K^\bullet/IK^\bullet, E^\bullet/I)
\end{align*}
and similarly for when dividing by $J$ or $I + J$, see
Derived Categories,
Lemma \ref{derived-lemma-morphisms-from-projective-complex}.
Derived Categories. Thus all the horizontal
arrows, except for possibly the middle one, are isomorphisms as the complexes
$K^\bullet/IK^\bullet$, $K^\bullet/JK^\bullet$, $K^\bullet/(I + J)K^\bullet$
are perfect complexes of $R/I$, $R/J$, $R/(I + J)$-modules, see
Proposition \ref{proposition-perfect-is-compact}.
It follows from the $5$-lemma (Homology, Lemma \ref{homology-lemma-five-lemma})
that the middle map is an isomorphism and the lemma follows by
Proposition \ref{proposition-perfect-is-compact}.
\end{proof}